\section{Planteamiento de la investigación}
\label{sec:planteamiento}

\subsection{Problema de investigación}
Las dimensiones originales no representan explícitamente todas las combinaciones geométricas de fibras y pivotes. Se desconoce si incorporar esas combinaciones mejora la estimación de la carga de primera falla y si la mejora se mantiene al excluir una familia del entrenamiento.

\noindent\textbf{Pregunta:} ¿En qué medida los momentos de sección y la razón logarítmica de volúmenes reducen el error de predicción frente a las entradas originales, manteniendo el mismo modelo y validación?

\subsection{Objetivos}
\noindent\textbf{General:} evaluar el aporte de características geométricas derivadas a la predicción de la carga de primera falla mediante una línea base y dos variantes de datos.
\begin{enumerate}
\item Caracterizar calidad, distribución y asociaciones de las variables para justificar su tratamiento.
\item Construir momentos de sección y razón de volúmenes sin utilizar respuestas posteriores al ensayo.
\item Comparar A, B y C mediante errores fuera de muestra por configuración y familia.
\end{enumerate}

\subsection{Hipótesis de trabajo}
\noindent\textbf{H1:} añadir los momentos (B) reduce el MAE por configuración al menos 5\% frente a la geometría original (A).

\noindent\textbf{H2:} añadir la razón logarítmica (C) reduce el MAE respecto de B.

Estas hipótesis organizan la evidencia ya disponible; no constituyen un preregistro ni pruebas de significación. Los criterios adicionales de adopción se especifican en la sección~\ref{sec:diseno}.

\subsection{Variables e indicadores}
\begin{table}[htbp]
\centering\small\setstretch{1.0}
\caption{Variables del estudio y forma de evaluación.}
\label{tab:variables-indicadores}
\begin{tabularx}{\linewidth}{@{}p{2.8cm}YY@{}}
\toprule
\rowcolor{TablaCabecera}
Rol & Variable o dimensión & Medición / indicador \\
\midrule
Entradas predictoras & Geometría de fibras, pivotes y malla & Dimensiones (mm), volúmenes (mm$^3$) y conteos; detalle en Tabla~\ref{tab:diccionario}. \\
\addlinespace
Salida objetivo & Carga de primera falla & Fuerza experimental en newtons (N); respuesta continua de regresión. \\
\addlinespace
Factor experimental & Representación de las entradas & A: original; B: añade momentos; C: añade la razón a B. \\
\addlinespace
Desempeño evaluado & Error de predicción y mejora incremental & MAE principal; RMSE, MedAE, $R^2$, reducción de MAE y tiempo como complementos. \\
\bottomrule
\end{tabularx}
\end{table}

\noindent Se controlan modelo, parámetros y particiones en A--B--C. La familia se usa para agrupar y evaluar generalización; no se interpreta como un efecto causal aislado.
\FloatBarrier
